\chapter{Wstęp}
Ostatnie lata przyniosły gwałtowny rozwój dużych modeli językowych, których zdolności w zakresie rozumienia i generowania języka naturalnego przekroczyły to, czego jeszcze dekadę wcześniej można się było spodziewać po systemach sztucznej inteligencji. Powszechna dostępność takich modeli w formie aplikacji konwersacyjnych sprawiła, że temat ten -- dotąd obecny głównie w środowisku naukowym -- stał się przedmiotem szerokiego zainteresowania opinii publicznej, przemysłu oraz badaczy z wielu dziedzin. Wzrost ten nie był przełomem naukowym sam w sobie, lecz katalizatorem: nasilił globalną ciekawość, w jaki sposób modele te faktycznie przetwarzają informacje i co dokładnie kryje się w ich wewnętrznej strukturze. Niniejsza praca dyplomowa stanowi próbę częściowego zaspokojenia tej ciekawości, przedstawiając proces badania tego zjawiska oraz możliwości wykorzystania zdobytej w ten sposób wiedzy do optymalizacji samych modeli.

\section{Motywacja}
\label{sec:motywacja}
Sztuczne sieci neuronowe (ang. \textit{Artificial Neural Networks}) stanowią dziś architekturę bazową dla większości modeli sztucznej inteligencji. Zbudowane z ogromnej liczby prostych jednostek obliczeniowych -- neuronów -- sieci te zawdzięczają swoją skuteczność zdolności do plastycznego dostosowywania się do stawianych przed nimi zadań w procesie treningu. Ta sama elastyczność jest jednak źródłem podstawowego ograniczenia: sieci neuronowe pozostają w praktyce czarnymi skrzynkami (ang. \textit{black box}) -- na wejściu dostarczane są dane, na wyjściu otrzymywany jest wynik, natomiast sposób, w jaki ten wynik powstaje, nie jest bezpośrednio i w pełni obiektywnie interpretowalny. Badacze mogą jedynie pośrednio wnioskować o skutkach aktywacji poszczególnych neuronów, stawiając hipotezy co do przyczyn ich pobudzenia. \newline
Zainteresowanie tym zjawiskiem nie jest nowe -- pytanie o to, w jaki sposób sieci neuronowe przetwarzają i reprezentują informacje, towarzyszy tej dziedzinie od dawna. Istotnie nasiliło się ono jednak wraz z pojawieniem się dużych modeli językowych, trenowanych na ogromnych zbiorach danych tekstowych i w rezultacie dysponujących dostępem do rozległej wiedzy -- w tym takiej, która nie powinna być powszechnie dostępna lub swobodnie wykorzystywana. Z tego powodu szczególnego znaczenia nabrało zrozumienie, w jaki sposób modele te \enquote{myślą} oraz w jaki sposób można świadomie kształtować ich zachowanie -- wzmacniając pożądane wzorce i osłabiając te niepożądane. Poddziedzina zajmująca się tego rodzaju analizą wewnętrznej struktury sieci neuronowych nazywana jest mechanistyczną interpretowalnością (ang. \textit{mechanistic interpretability}); choć wywodzi się ona z badań nad sieciami znacznie prostszymi niż współczesne modele językowe, to właśnie w kontekście tych ostatnich zyskała w ostatnich latach szczególne znaczenie praktyczne. \newline
Wiedza ta niesie ze sobą również potencjał praktyczny, wykraczający poza czysto poznawczy cel zrozumienia modelu, a powodów, dla których jest ona istotna, jest wiele. Z perspektywy bezpieczeństwa, wgląd w to, jakie wzorce zachowań model faktycznie wykorzystuje, pozwala wykryć te niepożądane, zanim ujawnią się one w praktycznym zastosowaniu. Z perspektywy praktycznej -- umożliwia lokalizowanie źródeł błędów czy niepożądanych korelacji wyuczonych przez sieć, a więc świadome debugowanie modelu, czego nie sposób osiągnąć, traktując go jako czarną skrzynkę. \newline
Trzecia motywacja, najbliższa tematyce niniejszej pracy, dotyczy optymalizacji: jeżeli da się wskazać, które komponenty sieci realizują istotne obliczenia, a które są względem zadania redundantne, otwiera to drogę do świadomego upraszczania modeli bez utraty ich funkcjonalności -- zarówno poprzez zmniejszenie ich zapotrzebowania na moc obliczeniową, jak i stworzenie mniejszych, wyspecjalizowanych wariantów dostosowanych do konkretnych zastosowań.

\section{Cele}
\label{sec:cele}
Pierwszym celem pracy jest lokalizacja i wizualizacja reprezentacji wiedzy faktycznej wyuczonej przez model językowy. Obiektem badań będzie mały, otwarto-źródłowy model językowy oparty na architekturze dekodera. Podobnie do podejść stosowanych w pracach \cite{templeton_scaling_monosemanticity} i \cite{lanza_extraction_analysis_multimodal_concepts}, do ekstrakcji cech reprezentacji wewnętrznych zastosowany zostanie Rzadki Autoenkoder (ang. \textit{Sparse Autoencoder} -- SAE), a konkretnie jego wariant Top-K \cite{gao_2024_scaling_evaluating_sparse_autoencoders}. Metoda ta pozwala uzyskać rzadkie, bardziej monosemantyczne rozkłady aktywacji, dając wgląd w to, jakie koncepty model \enquote{rozumie}, jak są one ze sobą powiązane i w jaki sposób grupuje on informacje w swojej wewnętrznej reprezentacji. Zebrane w ten sposób dane stanowią fundament dla kolejnych etapów pracy. \newline
Dysponując mapą konceptów wyuczonych przez model, zbadana zostanie możliwość wydzielenia z niego węższych modeli dziedzinowych za pomocą techniki przycinania (ang. \textit{pruning}). Modele takie, ograniczone do wybranej domeny wiedzy, miałyby zachowywać zbliżoną jakość w obrębie tej dziedziny przy jednoczesnej redukcji liczby parametrów i kosztów obliczeniowych. Efekt ten zostanie osiągnięty przez semantyczną klasyfikację wyekstrahowanych cech, a następnie usunięcie komponentów modelu nieistotnych dla wybranej domeny. \newline
Uzupełniająco zbadana zostanie możliwość innego wykorzystania wyznaczonych domen -- jako ekspertów w architekturze mieszaniny ekspertów (ang. \textit{Mixture of Experts}, MoE) \cite{jacobs_1991_adaptive_mixtures_of_local_experts}. W odróżnieniu od konwencjonalnych podejść do budowy modeli MoE, w których eksperci wyłaniani są w procesie treningu, w niniejszej pracy zbadana zostanie możliwość zdefiniowania ekspertów poprzez wydzielone, interpretowalne domeny wiedzy, wyznaczone na podstawie analizy cech SAE.

\section{Hipotezy i pytania badawcze}
\label{sec:hipotezy_pytania}
Realizacja celów przedstawionych w poprzedniej sekcji opiera się na czterech hipotezach badawczych. Każda z nich formułuje konkretne oczekiwanie co do własności reprezentacji wiedzy w modelu lub co do możliwości praktycznego ich wykorzystania; każdej towarzyszy odpowiadające pytanie badawcze wyznaczające zakres empirycznej weryfikacji.
\begin{enumerate}
    \item \textbf{Fakty i koncepty posiadają identyfikowalne reprezentacje w wewnętrznej strukturze modelu językowego.}
    Modele językowe nabywają wiedzę wyłącznie poprzez trening na tekście -- nie mają dostępu do żadnych zewnętrznych baz wiedzy ani słowników. Oznacza to, że informacje, które model jest w stanie poprawnie przywołać, muszą być zakodowane gdzieś w jego parametrach. Założenie o identyfikowalności tych reprezentacji jest zatem warunkiem wstępnym całego dalszego programu badawczego.
    \textit{Pytanie badawcze:} W jaki sposób można lokalizować i identyfikować reprezentacje wiedzy faktycznej w wewnętrznej strukturze modelu językowego?
    \newpage
    \item \textbf{Fakty powiązane semantycznie wykazują zbliżone lub sąsiadujące reprezentacje w przestrzeni wewnętrznej modelu.}
    Skoro model uczy się na podstawie wzajemnego współwystępowania słów i pojęć w tekście, fakty należące do tej samej dziedziny były mu wielokrotnie przedstawiane w podobnych kontekstach. Można zatem oczekiwać, że ich wewnętrzne reprezentacje będą geometrycznie zbliżone. Jest to zarazem warunek niezbędny do tego, by wiedza dziedzinowa dawała się sensownie wydzielić w kolejnych etapach pracy.
    \textit{Pytanie badawcze:} Czy fakty o podobnym znaczeniu wykazują wzajemne sąsiedztwo w przestrzeni wewnętrznej modelu i w jaki sposób można tę strukturę zmierzyć i zwizualizować?
    \item \textbf{Możliwe jest wydzielenie wąskiej dziedziny wiedzy z modelu poprzez technikę przycinania, z zachowaniem jakości w obrębie tej dziedziny.}
    Jeżeli wiedza faktyczna jest identyfikowalna i wykazuje przestrzenną spójność (hipotezy 1 i 2), to selektywne usuwanie komponentów modelu nieistotnych dla wybranej domeny powinno być możliwe bez poważnego pogorszenia jego wyników w zakresie tej domeny. Stopień, w jakim jest to faktycznie osiągalne, zależy od tego, na ile poszczególne dziedziny wiedzy są w modelu rozłączne..
    \textit{Pytanie badawcze:} Jakie są granice możliwości izolowania wiedzy dziedzinowej z modelu językowego poprzez technikę przycinania i jak mierzyć utratę jakości wynikającą z tego procesu?
    \item \textbf{Wyznaczone domeny wiedzy mogą stanowić podstawę do zdefiniowania ekspertów w architekturze MoE.}
    Jeżeli dziedziny wiedzy można wydzielić i opisać w sposób interpretowalny, stanowią naturalnych kandydatów na ekspertów w architekturze, w której przy każdym zapytaniu aktywuje się tylko ta część modelu kompetentna w danej dziedzinie. Takie podejście różni się od rozwiązań konwencjonalnych tym, że eksperci byliby zdefiniowani przez jawne kryteria semantyczne, a nie wyłaniani automatycznie w procesie treningu.
    \textit{Pytanie badawcze:} W jakim stopniu domeny wiedzy wyznaczone na podstawie analizy reprezentacji modelu mogą pełnić funkcję ekspertów w architekturze MoE?
\end{enumerate}